\begin{remarks}[7.6.7]
\label{IV.2.s7.6.7-fr}
\begin{enumerate}
\item[(i)] Las condiciones equivalentes de (7.6.6) significan también
que las fibras formales de $A$ son \emph{geométricamente regulares} (7.3.19, (iv)).
Ya se ha observado (\emph{loc. cit.}) que esta condición se satisface cuando
$A$ es \emph{un anillo de valoración discreta cuyo cuerpo de fracciones tiene
característica} 0.

\item[(ii)] El mismo razonamiento que en la primera parte de la
demostración de (7.6.4) muestra que, para un anillo semilocal noetheriano $A$,
las dos condiciones siguientes son equivalentes:
\begin{enumerate}
\item[a)] \emph{Para todo anillo cociente $B$ de $A$ que sea un dominio de
integridad, el completado $\widehat B$ es reducido.}
\item[a$'$)] \emph{Las fibras formales de $A$ son reducidas.}
\end{enumerate}

Teniendo en cuenta (\textbf 0, 23.1.7, (i)), estas dos condiciones implican la siguiente:
\begin{enumerate}
\item[b)] \emph{Para todo anillo cociente $B$ de $A$ que sea un dominio de
integridad, el cierre entero de $B$ es una $B$-álgebra finita.}
\end{enumerate}

Cuando $A$ es un anillo \emph{universalmente catenario}, puede demostrarse que,
recíprocamente, la condición $b)$ implica $a)$; no necesitaremos utilizar
este resultado.
\end{enumerate}
\end{remarks}
